\chapter{Experimental Program}
\label{ch:experiments-en}

The experimental program is described as the executable layer of the claim contract:
Stage~1/2 tests C01, Stage~3A tests C02, Stage~3B tests C03--C07, and QWake-FP C1/C2
tests C08--C11. Each block below separates the measurement objective, admissible unit
of inference, and transition condition to the next block.

The program was built sequentially. Each stage addressed a question left open by its
predecessor and did not replace earlier results. Stage~1/2, Stage~3A, Stage~3B, and
QWake therefore form one research chain: from quality and computational cost, through
internal mechanism, to an early compute-control decision under a registered
dangerous-accept criterion.

\input{generated/experimental_program_EN.tex}

\section{Stage 1: Initial Comparative Campaign}
\label{sec:stage1-experiment-en}

Stage~1 defined the initial empirical domain for RQ1 and the first part of C01. It used
FashionMNIST and MNIST, one \code{lenet\_classic} architecture, and four training
regimes: BP, Exact, FixedPred, and Strict. Ten independently trained models with
\code{model\_seed}=0--9 were executed for every dataset--method combination, producing
80 final cells.

PC hyperparameters were frozen before final execution: Exact used \(\eta=0.1\) and 20
state-inference steps; FixedPred used \(\eta=0.1\) and 10 steps; Strict used
\(\eta=0.05\) and 20 steps. BP had no corresponding internal-inference parameters.
Final-cell order was deterministically balanced so run order was not systematically
confounded with a method.

Test access was separated from pilot and validation stages. Final test evaluation was
performed only after predecessor-stage selection was frozen and only for a completed
run. Thesis-facing outputs are test accuracy, macro-F1, and registered computational
measures. The paired-analysis unit remained the independently trained model; examples
and batches were not treated as independent replications.

\section{Stage 2: Controlled Change of the Computational Path}
\label{sec:stage2-experiment-en}

Stage~2 completed C01 while retaining the 80-cell Stage~1 structure under a modified
implementation path. Its purpose was to separate properties of the studied regimes
from costs of a particular implementation. Datasets, architecture, methods, model
seeds, and final test protocol were retained, while the new campaign received its own
environment and artifact provenance chain.

This design permits matched cross-version comparison. Quality changes and
computational changes are analyzed separately. Zero quality change across versions is
not interpreted as universal implementation equivalence; it is bounded to the frozen
matrix, selected metrics, and finite model-seed set.

\section{Stage 3A: Layer-wise Diagnostics}
\label{sec:stage3a-experiment-en}

Stage~3A operationalized C02 and was restricted to FashionMNIST,
\code{lenet\_classic}, and ten independently trained models. Diagnostics were
extracted on validation data and did not reopen test access. BP, Exact, FixedPred, and
Strict were compared.

Gradient measures were cosine similarity, norm ratio, relative \(L_2\) error, and sign
agreement. Internal representations were compared using CKA and RSA. Metrics were
first aggregated within model and only then compared statistically at model-seed level.
Exact additionally served as a numerical control for gradient extraction and matching.

A separate depth analysis tested whether agreement with BP changed systematically from
early layers toward the output. Cross-layer CKA was retained descriptively and did not
turn layer pairs into independent statistical units.

\section{Stage 3B: Cost Profiling and Computational-domain Localization}
\label{sec:stage3b-experiment-en}

Stage~3B formed the C03--C07 surface and moved attention from final quality and
representation similarity to computational cost structure. The program comprised
several sequential tests.

\subsection{B0: Baseline Profiling}

B0 measured FixedPred and Strict in a frozen ROCm/float32 environment over a synthetic
scaling matrix. Three independent model seeds were used per configuration. Device and
CPU time, peak memory, and cost by execution region---including
\code{state\_inference}---were measured. B0 was not designed to prove structural
locality of the computational graph; it localized regions in which further
mechanistic cost attribution was meaningful.

\subsection{SI-MA0 and SI-MA1}

SI-MA0 tested the first registered \code{state\_inference} cost-accounting model. Its
negative outcome was preserved rather than used to retrospectively change a threshold.
SI-MA1 was a separate preregistered calibration, not a rerun of SI-MA0. It used ten
model seeds and 180 matched blocks. The primary estimand was a signed residual under a
predefined one-sided test. A negative signed residual means over-closure of the
accounting model, not negative physical runtime.

\subsection{B1 and B2: Equivalence Gates}

Two exact implementation candidates were then tested:
\code{isolated\_layer\_vjp} (B1) and \code{composite\_vjp} (B2). B1 compared the
candidate with the Stage~2 baseline over 120 registered pairs. B2 used 120 matched
triples and 240 pairwise comparisons among baseline, B1, and B2. Admission depended on
zero dangerous misses and passage of registered numerical, trajectory, structural,
observer, and provenance checks.

\subsection{Matched Profiling}

After EQ-B1 and EQ-B2, a separate matched-profiling matrix was executed. It contained
288 cells: 96 each for the Stage~2 baseline, \code{isolated\_layer\_vjp}, and
\code{composite\_vjp}. Cross-candidate correctness was checked on 96 matched blocks
outside the primary timing lane. Observer cost was not subtracted post hoc from the
primary measurements; observation-related structural measurements were retained
separately.

The matrix is engineering profiling, not 288 independent replications. A separate
descriptive analysis therefore applies the registered engineering-continuation screen
to four B1/B2-by-FixedPred/Strict groups of 16 configurations each. Decisions
\code{retain}, \code{conditional}, or \texttt{reject\_or\_revise} concern engineering
continuation only; the protocol forbids using the screen as a statistical or universal
superiority claim.

\section{QWake-FP: Sufficiency, Recognizability, and Economics}
\label{sec:qwake-experiment-en}

QWake-FP forms the C08--C11 experimental surface and continues the program after cost
localization. It separates three levels: whether preterminal records exist whose states
make the registered early action task-admissible; whether a non-zero subset can be
recognized before action with zero observed dangerous accepts on the calibration
surface; and whether the decision remains beneficial after full decision-cost
accounting.

The registered QWake-FP early action is \texttt{ANALYTIC\_COMPLETION} with identifier
\nolinkurl{fixedpred_eta1_wavefront_completion_v1}. It replaces the remaining
canonical iterative sweeps with a bounded analytic completion. The exact path
\nolinkurl{complete_suffix_stage2_baseline_v1} executes the full remaining suffix and
serves as reference and fallback. The \texttt{EARLY\_ADMISSIBLE} label is created only
after comparing required responses of those two branches; it does not mean that no
computation occurs after the decision epoch. Measured analytic-completion cost is part
of the registered C2 compute-cost component.

C1 produced sealed calibration trajectories and post-action sufficiency labels. The
exact oracle was used only for later evaluation and was not a C2 pre-action feature.
The frozen C1 request uses three model seeds and 36 batches, giving 108 trajectories.
The contract includes decision epochs \(t\in[0,K_{\mathrm{ref}}]\); with
\(K_{\mathrm{ref}}=6\), each trajectory yields seven records and the complete C2
surface contains 756 records: 648 preterminal records at \(t<6\) and 108 terminal-
boundary records at \(t=6\). C2 evaluates the whole surface and a preregistered finite
family of 2,625 scalar rules. The generator is restricted to selected scalar features
and one-dimensional threshold rules; C2 is therefore a bounded search, not an
exhaustive search over all recognizers.

For every rule, coverage, dangerous-accept count, and aggregate net saving under frozen
full decision-cost accounting are recorded. Eligibility requires zero dangerous
accepts, non-zero coverage, and positive aggregate net saving simultaneously. If no
rule satisfies all three, C2 terminates with a bounded negative result and the protocol
does not open C3.

\section{Experimental-program Stopping Principle}
\label{sec:experiment-stopping-en}

The program stops at registered decision boundaries rather than at a desired sign of
the result. Stage~3B preserves the negative SI-MA0 result before opening a narrower
observer-cost calibration. B1/B2 equivalence gates precede matched profiling, and the
engineering-continuation screen remains separate from equivalence. QWake-FP similarly
requires a C2 rule-freeze receipt before C3 can open. Because C2 found no eligible rule,
no rule-freeze receipt was created and C3 did not open. In the claim contract, C10 and
C11 therefore remain \texttt{not\_tested}: C10 requires a separate measurement of the
marginal execution cost of a minimal recognizer, while C11 requires confirmatory C3 to
have actually opened and independently confirmed a selected C2 rule. No post-hoc
threshold extension, feature addition, or cost redefinition is used to alter that
terminal outcome.
